%**********************************************************************
%* Niets toevoegen aan dit bestand zonder toestemming van uw promotor *
%**********************************************************************
\documentclass[12pt,a4paper,oneside]{book}

%%%%%%%%%%%%%%%%%%%%%%%%%%%%%%%%%%%%%%%%%%%%%%%%%%%%%%%%%%%%%%%%%%%%%%%%%%%%%%%%%%
%% Volgende regel activeren voor eReader output:
%\usepackage[papersize={130mm,180mm},margin=2mm]{geometry}

%% Volgende regel activeren voor A4 output (bovenstaande terug afzetten)
\usepackage{a4wide} % Iets meer tekst op een bladzijde
%%%%%%%%%%%%%%%%%%%%%%%%%%%%%%%%%%%%%%%%%%%%%%%%%%%%%%%%%%%%%%%%%%%%%%%%%%%%%%%%%%

%% Packages %%
\usepackage{a4wide}                     % Iets meer tekst op een bladzijde
\usepackage[dutch]{babel}               % Voor nederlandstalige woordsplitsing (en chapter hoofdingen)
\usepackage{amsmath}                    % Uitgebreide wiskundige mogelijkheden
\usepackage{url}                        % Om urls te verwerken
\usepackage{graphicx}                   % Om figuren te kunnen verwerken
\usepackage[latin1]{inputenc}           % Om niet ascii karakters te kunnen typen
\usepackage{listings}                   % improve the code environments
%\usepackage{lmodern}
%\usepackage[T1]{fontenc}
\usepackage{textcomp}
\usepackage{pdfpages}                   % Om pdf documenten in te kunnen voegen (titelpagina)
\usepackage{lipsum}                     % Om lorem ipsum in te voegen
\usepackage{sty/eukdate}                % Om de datum in te voegen 
\usepackage[DoggensLenny]{sty/fncychap} % Added for nicer formatting: (kaders rond titels)
\usepackage{fix-cm}
\usepackage{parskip}                    % Geen inspringing aan het begin van een alinea
%\usepackage{sectsty}
\usepackage{mdwlist}
\usepackage{sty/eukdate}                % Om data in te voegen
\usepackage{caption}
\usepackage{float}                      % Om afbeelding geforceerd op een plek te plaatsten met 'H' parameter
\restylefloat{figure}
\usepackage{multirow}                   % Weergave van complexere tabellen
\usepackage[scaled]{sty/inconsolata}    % Instelling lettertype
\usepackage{verbatim}

\usepackage[pagebackref=true]{hyperref}
\hypersetup{                            % Markup for links in text, final versie enkel met zwarte links!!  %%
    colorlinks = false,                 %false=rode kaders rond tekst, true=links in de tekst zelf in kleur
    linkcolor = black,
    anchorcolor = black,
    citecolor = blue,
    filecolor = black,
    urlcolor = black
}

\renewcommand*\familydefault{\sfdefault}                        % Only if the base font of the document is to be sans serif
%\allsectionsfont{\usefont{OT1}{phv}{bc}{n}\selectfont}
  
\hyphenpenalty=5000                                             % Minder snel woorden splitsen
  \tolerance=1000

%% Custom commands %%
\newcommand{\npar}{\par \vspace{2.3ex plus 0.3ex minus 0.3ex}}  % Enkele mm vertikale ruimte openlaten
\newcommand{\ntpar}{\par \vspace{1.3ex plus 0.3ex minus 0.3ex}}
\newcommand{\nhpar}{\par \vspace{4.3ex plus 0.3ex minus 0.3ex}}

\newcommand{\HRule}{\rule{\linewidth}{0.5mm}}                   % Horizontale lijn

\setcounter{tocdepth}{2}                                        % weergeven tot subsection in table of contents

% \usepackage[final]{sty/TODO}           %TODO package (build error als er nog een todo in het document staat)
%\usepackage[silent]{sty/TODO}           %TODO package (geen output, om een tussentijdse draft zonder de 'todos' te maken)
\usepackage{sty/TODO}                    %TODO package (standaard modus)


\usepackage{eso-pic}
\newcommand\BackgroundPic{%
\put(0,0){%
\parbox[b][\paperheight]{\paperwidth}{%
\vfill
\centering
% \includegraphics[width=\paperwidth,height=\paperheight]{figures/titelpagina_ap_lowestres.jpg}
% \includegraphics[width=\paperwidth,height=\paperheight]{figures/titelpagina_ap_lowres.jpg}
\includegraphics[width=\paperwidth,height=\paperheight]{figures/titelpagina_ap_mediumres.jpg}

\vfill
}}\hfill
}

\newcommand\BackgroundPicAP{%
\put(220,-372){%
\parbox[b][\paperheight]{\paperwidth}{%
\vfill
\centering
\includegraphics[width=0.44\paperwidth,height=0.075\paperheight,keepaspectratio]{figures/AP.png}
\vfill
}}
\hfill
}

%*************************************************************************************
%* Plaats eventuele extra toegevoegde packages hier (na toestemming van uw promotor) *
%*************************************************************************************



